\subsection{Estrategia progresiva y autoridad de registro}
Synaptix desplegará por proceso y zona. E1 priorizará torre y zarpes, inventario y autorización de puestos; después incorporará Escuela náutica, combustible y medición. La primera ola responde a la custodia y al estado de embarcaciones utilizados por los demás procesos. Escuela dependerá de identidad, autorizaciones y retiros verificados \parencite[cap.~17, secc.~17.6]{caso07}.

Las olas E1 se prepararán antes del mes 13 y se habilitarán entre el 1 y el 14 de diciembre de 2027, dentro del mes 13 del escenario seleccionado; desde el 15 no habrá habilitaciones nuevas hasta terminar el congelamiento. El alcance completo permanecerá en observación durante las cuatro semanas finales. E2 incorporará progresivamente varadero, contratistas, administración y gestión ambiental completa durante el 19, para observar el conjunto durante las cuatro semanas finales del 20. Las innovaciones conservarán sus pilotos de SD13; invernada no se forzará a meses 19--20.

Cada dato tendrá una autoridad de escritura por período. Antes del cambio administrativo el legado mantendrá su autoridad y las cargas controladas permitirán comprobar la nueva solución. Los hechos del recinto conservarán autoridad local según SD4/SD5. No se impondrá digitación ordinaria doble; se compararán extracciones y registros controlados. Cada ola tendrá autorizador, habilitación y procedimiento de suspensión y retorno.

\subsection{Etapa 1: marcha blanca 13--15 y producción 16}
La primera etapa habilitará seguridad y continuidad del recinto, escuela, inventario, combustible y medición individual, conforme a SD3. La preparación incluirá ensayos de migración, pruebas de continuidad y personas capacitadas. Durante marcha blanca se mantendrán los controles vigentes de la marina y se conciliará diariamente, sin sustituir una observación segura por una señal sensórica. La recepción verificará CA-01--CA-08, CA-10, CA-13 y CA-23 y la medición inicial de CA-18; CA-19 conserva su fecha ambiental externa. El paso oficial al mes 16 requerirá cierre aceptado y activará cuatro semanas de estabilización.

\subsection{Etapa 2: marcha blanca 19--20 y producción 21}
La segunda etapa completará varadero, contratistas, ambiente y administración, manteniendo E1 en producción. Sus ensayos comprobarán versiones e interfaces de E2 y protección de los datos compartidos. La recepción verificará CA-09, CA-11, CA-12, CA-14--CA-22 y CA-24 y conservará los resultados E1. El paso del mes 21 iniciará simultáneamente los 36 meses de Operación de ambos alcances; la estabilización inicial se imputará a ese servicio. El cambio administrativo respetará la consulta independiente y la autoridad definida en SD5.

\subsection{Condiciones de avance y cierre de marcha blanca}
El avance exigirá recepción del alcance previo, preparación de usuarios, pruebas aceptadas, integridad y retorno disponible. El cierre cumplirá conjuntamente: ningún incidente crítico ni alto abierto atribuible a la solución; volumen real comprometido durante las cuatro semanas finales; disponibilidad y tiempos sostenidos durante ese período; conciliación sin diferencias no explicadas; capacitación y certificación exigida; y acta suscrita por la Contraparte Técnica \parencite[art.~17.3, p.~13]{bases-admin}.

Se compromete cobertura del 100\,\% de hechos reales exigibles de los procesos habilitados durante esas cuatro semanas, con recuentos diarios contrastados. La cobertura no equivale al peak: se registrarán cantidades absolutas, mezcla de operaciones y períodos sin actividad. Calendario estacional y fecha de inicio determinarán el volumen absoluto a aprobar antes de marcha blanca. Operaciones inventadas y carga sintética no acreditarán volumen real.

Los indicadores diarios incluirán disponibilidad, tiempos de flujos de SD3, operaciones conciliadas, diferencias por conjunto y moneda, duplicados, autorizaciones y usuarios preparados. Los umbrales conservarán los criterios CA de cada etapa y los compromisos de SD4/SD5. T-17 recogerá evidencia y observaciones resueltas conforme al artículo 18. Si falta una condición al cierre, Synaptix extenderá la marcha blanca a su costo, sin cargo adicional ni desplazamiento de fases siguientes y bajo las condiciones del artículo 80 \parencite[arts.~17.3 y 18, p.~13]{bases-admin}.

\subsection{Calendario protegido, marejadas y retorno}
No se intervendrán los sistemas entre el 15 de diciembre y el 15 de marzo ni durante las 14 fechas náuticas; se evitará la capacitación en el peak protegido. Las campañas de varada de abril--mayo y septiembre--noviembre bloquearán toda intervención que afecte al varadero. Se protegerán además maniobras y clases; el campo de boyas tiene acceso sólo por mar. Se reservan inicialmente dos ventanas alternativas por intervención marina y diez días hábiles por campaña de despliegue en terreno, antes de su hito. Son hipótesis por contrastar con calendario y cancelaciones, sin modificar meses 16/21 ni avisos de suspensión de 24--48 horas.

El aviso de marejada suspenderá la intervención y activará la alternativa con coordinación segura. La reserva se distinguirá de holguras libre y total de CPM y del retorno. La red y la sensibilidad con recursos incluidas en este formulario cuantifican el encaje bajo el escenario seleccionado. Las fechas efectivas de eventos, campañas y autorizaciones se comprobarán antes de ejecutar; el cálculo no acredita una instalación, una prueba ni un acta futura \parencite[caps.~10, 13 y 17, secc.~17.5]{caso07}.

El corte administrativo conserva 22:00--02:00, decisión antes de 00:30 y 90 minutos exclusivos de retorno previo. El retorno posterior conserva la estimación separada de 120 minutos, por comprobar con el mayor volumen previsto de cambios y respuestas pendientes. Si el legado no representa un hecho, se mantendrá el diario conciliado y la corrección controlada, sin eliminar operaciones ni repetir pagos.

\subsection{Pruebas y preparación de cada etapa}
La aceptación incluirá migración, usuario, integración, regresión, seguridad ofensiva, accesibilidad WCAG 2.2 AA, resiliencia y recuperación. La carga comprobará los umbrales a 1,5 veces el peak declarado en PREPROD. El estrés aumentará progresivamente la demanda hasta identificar el punto de quiebre, en un entorno aislado; el informe conservará curva de respuesta, saturación, recursos, degradación y recuperación durante y después del peak; componentes conservarán cobertura mínima de 70\,\% de lógica de negocio y ninguna prueba en falla. No se cerrará una etapa con hallazgos críticos o altos abiertos \parencite[cap.~20, secc.~20.1]{bases-transversales}.

La capacitación se organizará por perfil y turno, con ejercicios y evaluación antes de las ventanas protegidas. Incluirá personal estacional y monitores nuevos, sin suponer un área TI del cliente. La estabilización posterior tendrá presencia en pantalanes y varadero, incluidos fines de semana pertinentes; la dotación y los turnos seleccionados se detallan en el dimensionamiento y cuadrantes de T-15. El esfuerzo de seguimiento de datos no acredita esa cobertura física completa \parencite[cap.~22, RT-22.01--RT-22.07]{bases-transversales}.
